\chapter{Complexity — DAG Hard Functions}
%%% generated by the blueprint pipeline from TCSlib.Complexity.CircuitComplexity.DAGHardFunctions
\section{Overview}

Declarations of the Lean module \texttt{TCSlib.Complexity.CircuitComplexity.DAGHardFunctions}.

\section{Declarations}

\begin{definition}[Encoding a gate]
\lean{BoolCircuit.GateCode}
\label{BoolCircuit.GateCode}
\leanok
For a natural number $S$, a gate code consists of a choice among three labels, together
with two optional indices, each either absent or an element of $\{0,\ldots,S-1\}$.
\end{definition}

\begin{definition}[Encoding a circuit]
\lean{BoolCircuit.CircCode}
\label{BoolCircuit.CircCode}
\leanok
\uses{BoolCircuit.GateCode}
For a natural number $S$, a circuit code consists of a number from $0$ through $S$, a
function assigning a gate code to each of the $S$ gate positions, and a designated
output index from $0$ through $S-1$.
\end{definition}

\begin{definition}[Decoding gate kinds]
\lean{BoolCircuit.decodeKind}
\label{BoolCircuit.decodeKind}
\leanok
\uses{BoolCircuit.GateKind}
Decode the three labels in order as AND, OR, and NOT: the first label represents AND,
the second represents OR, and the third represents NOT.
\end{definition}

\begin{definition}[Gate-kind encoding]
\lean{BoolCircuit.encodeKind}
\label{BoolCircuit.encodeKind}
\leanok
\uses{BoolCircuit.GateKind}
The three gate kinds are encoded by the elements $0$, $1$, and $2$ of $	ext{Fin}(3)$:
conjunction maps to $0$, disjunction to $1$, and negation to $2$.
\end{definition}

\begin{theorem}[Decoding an encoded gate kind]
\lean{BoolCircuit.decodeKind_encodeKind}
\label{BoolCircuit.decodeKind_encodeKind}
\leanok
\uses{BoolCircuit.GateKind, BoolCircuit.decodeKind, BoolCircuit.encodeKind}
\difficulty{1}
For each gate kind $k \in \{\mathrm{AND},\mathrm{OR},\mathrm{NOT}\}$, decoding its
encoding returns $k$.
\end{theorem}

\begin{definition}[Gate decoding]
\lean{BoolCircuit.decodeGate}
\label{BoolCircuit.decodeGate}
\leanok
\uses{BoolCircuit.DAGGate, BoolCircuit.GateCode, BoolCircuit.decodeKind}
For a natural number $S$, a gate code consists of a kind code in $	ext{Fin}(3)$ and two
optional argument indices in $	ext{Fin}(S)$. Decode the kind code as conjunction,
disjunction, or negation according as it is $0$, $1$, or $2$, and form the gate's
argument list by taking the values of the two present indices in order, omitting absent
ones.
\end{definition}

\begin{definition}[Argument encoding]
\lean{BoolCircuit.encodeArg}
\label{BoolCircuit.encodeArg}
\leanok
Given a natural number $S$ and an optional natural-number argument $a$, encode it as an
optional element of $	ext{Fin}(S)$: an absent argument, or a present $a
ot< S$, is encoded as absent; a present $a<S$ is encoded by its corresponding element of
$	ext{Fin}(S)$.
\end{definition}

\begin{definition}[Gate encoding]
\lean{BoolCircuit.encodeGate}
\label{BoolCircuit.encodeGate}
\leanok
\uses{BoolCircuit.DAGGate, BoolCircuit.GateCode, BoolCircuit.encodeArg, BoolCircuit.encodeKind}
For a natural number $S$, encode a gate by the triple consisting of its kind code and
the encodings of its first and second arguments, where arguments beyond the first two
are ignored and a missing argument is treated as absent.
\end{definition}

\begin{theorem}[Gate encoding and decoding]
\lean{BoolCircuit.decodeGate_encodeGate}
\label{BoolCircuit.decodeGate_encodeGate}
\leanok
\uses{BoolCircuit.DAGGate, BoolCircuit.decodeGate, BoolCircuit.decodeKind_encodeKind, BoolCircuit.encodeArg, BoolCircuit.encodeGate}
\difficulty{4}
Let $S$ be a natural number and let $g$ be a gate with at most two arguments, each a
natural number less than $S$. Encoding $g$ using the three gate labels and two bounded
argument slots, then decoding that code, recovers $g$ exactly.
\end{theorem}

\begin{definition}[Gates represented by a circuit code]
\lean{BoolCircuit.gatesOf}
\label{BoolCircuit.gatesOf}
\leanok
\uses{BoolCircuit.CircCode, BoolCircuit.DAGGate, BoolCircuit.decodeGate}
For a circuit code with capacity $S$ and a specified number of gates, take the list of
that many gates by decoding the corresponding gate codes at indices $0,1,\ldots$ in
order.
\end{definition}

\begin{definition}[Evaluation of a circuit code]
\lean{BoolCircuit.evalCode}
\label{BoolCircuit.evalCode}
\leanok
\uses{BoolCircuit.CircCode, BoolCircuit.DAGGate.eval, BoolCircuit.gatesOf, BoolCircuit.runWith}
For a natural number $n$, a circuit code of capacity $S$, and an input
$x:	ext{Fin}(n)	o	ext{Bool}$, evaluate the decoded gate list on the list of input
values. Return the resulting value at the circuit's output index if present, and return
false if that index has no value.
\end{definition}

\begin{theorem}[Encoding a bounded fan-in circuit]
\lean{BoolCircuit.exists_code}
\label{BoolCircuit.exists_code}
\leanok
\uses{BoolCircuit.CircCode, BoolCircuit.DAGCircuit, BoolCircuit.DAGCircuit.IsFaninTwo, BoolCircuit.DAGCircuit.eval, BoolCircuit.DAGCircuit.size, BoolCircuit.GateCode, BoolCircuit.decodeGate_encodeGate, BoolCircuit.encodeGate, BoolCircuit.evalCode, BoolCircuit.gatesOf}
\difficulty{4}
Let $n$ and $S$ be nonnegative integers, and let $C$ be a directed acyclic Boolean
circuit on $n$ input variables, with its gates in topological order and a valid output
vertex. Suppose each gate has at most two inputs, its input indices are distinct, and
every NOT gate has exactly one input. If the circuit size, counting the input vertices
and gates, is at most $S$, then there is a code of size parameter $S$ whose evaluation
agrees with $C$ on every input assignment.
\end{theorem}

\begin{definition}[Functions computed by bounded DAG circuits]
\lean{BoolCircuit.computableDAG}
\statementsource{Vad12}{
  pdf-bf2ce5422088-p029-b001,
  pdf-bf2ce5422088-p217-b003
}
\label{BoolCircuit.computableDAG}
\leanok
\uses{BoolCircuit.DAGCircuit, BoolCircuit.DAGCircuit.IsFaninTwo, BoolCircuit.DAGCircuit.eval, BoolCircuit.DAGCircuit.size}
For natural numbers $n$ and $S$, the set of functions from $	ext{Fin}(n)	o	ext{Bool}$ to
$	ext{Bool}$ computed by an acyclic gate circuit whose gates each have at most two
inputs and whose size is at most $S$.
\end{definition}

\begin{theorem}[Count of functions computed by small DAG circuits]
\lean{BoolCircuit.card_computable_dag_le}
\label{BoolCircuit.card_computable_dag_le}
\leanok
\uses{BoolCircuit.CircCode, BoolCircuit.computableDAG, BoolCircuit.evalCode, BoolCircuit.exists_code}
\difficulty{3}
For natural numbers $n$ and $S$, the number of Boolean functions on $n$ variables
computed by well-formed fan-in-two directed acyclic circuits of size at most $S$ is at
most $(S+1)\bigl(3(S+1)^2\bigr)^S S$.
\end{theorem}

\begin{theorem}[Existence of a hard function under a circuit-count bound]
\lean{BoolCircuit.exists_not_eval_dag_of_lt}
\label{BoolCircuit.exists_not_eval_dag_of_lt}
\leanok
\uses{BoolCircuit.DAGCircuit, BoolCircuit.DAGCircuit.IsFaninTwo, BoolCircuit.DAGCircuit.eval, BoolCircuit.DAGCircuit.size, BoolCircuit.card_computable_dag_le, BoolCircuit.computableDAG}
\difficulty{4}
Let $n,S$ be natural numbers satisfying $(S+1)\bigl(3(S+1)^2\bigr)^S S<2^{2^n}$. Then
there is a Boolean function on $n$ input bits such that every well-formed acyclic
Boolean circuit with at most two inputs to each gate and size at most $S$ disagrees with
it on at least one input.
\end{theorem}

\begin{theorem}[Circuit-counting inequality]
\lean{BoolCircuit.count_lt}
\label{BoolCircuit.count_lt}
\leanok
\difficulty{5}
Let $n,S$ be natural numbers with $1 \le n$ and $10nS \le 2^n$. Then
$(S+1)\bigl(3(S+1)^2\bigr)^S S < 2^{2^n}$.
\end{theorem}

\begin{theorem}[Existence of a hard Boolean function for DAG circuits]
\lean{BoolCircuit.exists_hard_function_dag}
\label{BoolCircuit.exists_hard_function_dag}
\leanok
\uses{BoolCircuit.DAGCircuit, BoolCircuit.DAGCircuit.IsFaninTwo, BoolCircuit.DAGCircuit.eval, BoolCircuit.DAGCircuit.size, BoolCircuit.count_lt, BoolCircuit.exists_not_eval_dag_of_lt}
\difficulty{1}
For every natural number $n>1$, there is a Boolean function on $n$ input bits such that
every well-formed acyclic Boolean circuit with at most two inputs to each gate and size
at most $2^n/(10n)$ disagrees with the function on at least one input.
\end{theorem}

\begin{theorem}[A language hard for bounded-size DAG circuits]
\lean{BoolCircuit.exists_language_hard_dag}
\label{BoolCircuit.exists_language_hard_dag}
\leanok
\uses{BoolCircuit.DAGCircuit.eval, BoolCircuit.exists_hard_function_dag, Language.InSIZE}
\difficulty{5}
There is a language $L$ over Boolean strings such that, for every function
$T\colon\mathbb{N}\to\mathbb{N}$ for which some $n$ satisfies $1<n$ and $T(n)\le
2^n/(10n)$, no family of well-formed acyclic Boolean circuits with at most two inputs to
each gate, size at most $T(n)$ on inputs of length $n$, recognizes $L$.
\end{theorem}
